\begin{ouverture}[ensembles]{16}{Raisonner et faire la synthèse}{Partie 5 · Méthodes transversales}
{Ensembles · Logique · Démonstrations · Problèmes mixtes}
\objectif{Lire et écrire une proposition avec ses hypothèses.}
\objectif{Choisir un contre-exemple ou une méthode de preuve.}
\objectif{Mobiliser plusieurs chapitres dans un même problème.}
\prerequis{Ces ressources se travaillent progressivement toute l’année.}
\end{ouverture}
\pagestyle{otMain}\essentiel\fontsize{11}{15}\selectfont
\essSection{1}{Ensembles et propositions}
\begin{definition}$x\in E$ signifie que $x$ est un élément de $E$ ; $A\subset E$ signifie que tout élément de $A$ appartient à $E$ : $A$ est un sous-ensemble de $E$. L’ensemble vide est $\varnothing$. Pour $A,B\subset E$ : l’intersection $A\cap B$ contient les éléments communs à $A$ et $B$ ; la réunion $A\cup B$ ceux qui sont dans $A$ ou dans $B$ ; le complémentaire de $A$ dans $E$, noté $\overline A$ ou $E\setminus A$, ceux de $E$ qui ne sont pas dans $A$. Le cardinal $\operatorname{Card}(A)$ d’un ensemble fini est son nombre d’éléments. Le produit cartésien $A\times B$ est l’ensemble des couples $(a;b)$ avec $a\in A$ et $b\in B$.\end{definition}
\begin{definition}Une proposition est un énoncé mathématique qui est soit vrai, soit faux. Elle peut contenir une variable : « $n$ est pair » est vraie pour $n=4$ et fausse pour $n=7$.\end{definition}
\begin{propriete}« Et » exige les deux conditions ; « ou » signifie au moins l’une des deux, y compris les deux. Nier $x<3$ donne $x\ge3$. Nier « $x$ est pair et positif » donne « $x$ n’est pas pair ou n’est pas positif ».\end{propriete}
\begin{definition}Une implication « si $P$, alors $Q$ » affirme que l’hypothèse $P$ entraîne $Q$. Sa réciproque échange $P$ et $Q$. Sa contraposée est « si non $Q$, alors non $P$ » ; elle a la même valeur de vérité que l’implication. Une équivalence $P\iff Q$, lue « $P$ si et seulement si $Q$ », exige les deux implications. Par exemple, pour un réel $x$, $x^2=4\iff x=-2\text{ ou }x=2$.\end{definition}
\begin{attention}Une implication n’autorise pas automatiquement sa réciproque. « Multiple de 4 entraîne pair » est vraie, mais « pair entraîne multiple de 4 » est fausse.\end{attention}
\clearpage
\essSection{2}{Chercher, démontrer et communiquer}
\begin{savoirfaire}[Une preuve adaptée]
Pour une affirmation universelle, raisonner sur un objet quelconque avec les hypothèses. Pour une existence, fournir un exemple qui convient. Pour réfuter une affirmation universelle, un seul contre-exemple suffit. Les formulations quantifiées sont écrites en mots ; les symboles de quantification ne sont pas requis en Seconde.
\end{savoirfaire}
\begin{savoirfaire}[Trois démarches]
Disjonction : étudier tous les cas possibles. Contraposée : démontrer non $Q$ entraîne non $P$. Absurde : supposer la conclusion fausse puis obtenir une contradiction avec les hypothèses ou une propriété établie.
\end{savoirfaire}
\begin{exemple}Pour montrer qu’un entier dont le carré est pair est pair, on peut utiliser la contraposée : le carré d’un impair est impair. La preuve générale se trouve au chapitre 1 ; un simple test numérique ne suffirait pas.\end{exemple}
\otIntertitre{Une rédaction complète}
Nommer les objets ; annoncer les conditions ; présenter les étapes dans l’ordre ; conclure dans le contexte et avec les unités. Un schéma sert à chercher et expliquer, mais une propriété lue à l’œil doit être justifiée.
\par\medskip Les exercices 11 à 15 mélangent les domaines. Rechercher d’abord les liens utiles entre leurs données plutôt que tenter d’appliquer le dernier cours appris.
\clearpage\rubrique{Découverte}
\begin{blocLivre}{ex-16-1}
\begin{exercice}[1]{Appartenance et inclusion}
\begin{questions}
\item Pour $A=\{1;2;4\}$, dire si $2\in A$, $\{2\}\subset A$, $3\in A$ et $\varnothing\subset A$ sont vraies.
\item Dans $E=\{1;2;3;4;5\}$, on donne aussi $B=\{2;3;5\}$. Déterminer $A\cap B$, $A\cup B$, $\overline A=E\setminus A$ et $\operatorname{Card}(A\cup B)$.
\end{questions}
\end{exercice}
\end{blocLivre}
\begin{blocLivre}{ex-16-2}
\begin{exercice}[1]{Un produit cartésien}
Pour $A=\{0;1\}$ et $B=\{a;b;c\}$, lister $A\times B$ et donner son cardinal.
\end{exercice}
\end{blocLivre}
\begin{blocLivre}{ex-16-3}
\begin{exercice}[1]{Nier une condition}
Formuler la négation de « $x<5$ », puis de « $x\ge2$ et $x\le7$ ».
\end{exercice}
\end{blocLivre}
\begin{blocLivre}{ex-16-4}
\begin{exercice}[1]{Implication et réciproque}
Écrire la réciproque de « si un entier est multiple de 6, il est pair ». Donner la valeur de vérité de chacune des deux phrases.
\end{exercice}
\end{blocLivre}
\begin{blocLivre}{ex-16-5}
\begin{exercice}[1]{Prouver une existence}
Existe-t-il un réel $x$ tel que $x^2=2x$ ? Donner toutes les solutions. Quelle différence y a-t-il entre « il existe » et « pour tout » ?
\end{exercice}
\end{blocLivre}

\clearpage\rubrique{Entraînement}
\begin{blocLivre}{ex-16-6}
\begin{exercice}[2]{Un contre-exemple décisif}
Réfuter : « Pour tout réel $x$, si $x^2>4$, alors $x>2$. » Proposer une conclusion correcte.
\end{exercice}
\end{blocLivre}
\begin{blocLivre}{ex-16-7}
\begin{exercice}[2]{Une disjonction de cas}
Démontrer que le produit de deux entiers consécutifs est pair en distinguant les deux parités possibles du premier.
\end{exercice}
\end{blocLivre}
\begin{blocLivre}{ex-16-8}
\begin{exercice}[2]{Une preuve par l’absurde}
Démontrer qu’il n’existe pas de plus grand entier naturel.
\end{exercice}
\end{blocLivre}
\begin{blocLivre}{ex-16-9}
\begin{exercice}[2]{Une contraposée}
Pour un entier $n$, démontrer : « si $n^2$ est impair, alors $n$ est impair », en utilisant la contraposée.
\end{exercice}
\end{blocLivre}
\begin{blocLivre}{ex-16-10}
\begin{exercice}[2]{Le statut d’un dessin}
Un logiciel affiche trois points presque alignés. Quelles données permettent de démontrer l’alignement ? Proposer un exemple de trois points très proches d’une même droite mais non alignés.
\end{exercice}
\end{blocLivre}

\clearpage\rubrique{Problèmes}
\begin{blocLivre}{ex-16-11}
\begin{exercice}[3]{Synthèse : une aire maximale}
Un rectangle a un périmètre de 28 cm. Exprimer son aire en fonction d’un côté $x$, déterminer son maximum et les dimensions correspondantes.
\label{fig-ex-16-11}
\figureGeo{synthese-rectangle}
\end{exercice}
\end{blocLivre}
\begin{blocLivre}{ex-16-12}
\begin{exercice}[3]{Synthèse : contrôler une pièce}
Une fabrication de 500 pièces compte 40 défectueuses. Un contrôle donne 36 signaux positifs parmi les défectueuses et 23 parmi les conformes. Calculer la sensibilité et la part de défauts parmi les signaux positifs.
\end{exercice}
\end{blocLivre}
\begin{blocLivre}{ex-16-13}
\begin{exercice}[3]{Synthèse : coûts et seuil}
Une sortie coûte 180 euros de transport et 12 euros par participant. Chaque participant paie 18 euros. Exprimer le résultat financier pour $n$ participants et déterminer le seuil de rentabilité.
\end{exercice}
\end{blocLivre}
\begin{blocLivre}{ex-16-14}
\begin{exercice}[3]{Synthèse : un lieu et une droite}
$A(0;0)$ et $B(4;0)$. Trouver le point de la droite $y=x+1$ qui est équidistant de $A$ et de $B$.
\label{fig-ex-16-14}
\figureGeo{synthese-mediatrice}
\end{exercice}
\end{blocLivre}
\begin{blocLivre}{ex-16-15}
\begin{exercice}[3]{Synthèse : juger une conclusion}
Une série de 20 mesures a une moyenne de 10. Un élève affirme que 10 de ces mesures sont forcément inférieures ou égales à 10. Réfuter sa conclusion avec une série explicite, puis préciser l’indicateur qui parle de partage en deux.
\end{exercice}
\end{blocLivre}

\clearpage
